%% annex/p04/decision-problems-and-arithmetical-hierarchy.tex
%% Parallel proof fragment for the v1.1 lines 8248--8897 migration.
%% Complete proofs are compressed once toward the 40--60 percent target;
%% sketch omissions remain explicit.

\begin{annexitem}{annex:p04-hierarchy.1}{thm:p04-fol-validity-ce}{Enumerability without decidability}
\begin{proof}
Enumerate empty-premise finite derivations with sentence conclusions.
Soundness makes every output valid; completeness makes every valid sentence
appear, so validity is computably enumerable.  A decision procedure would
contradict Theorem~\ref{thm:p04-fol-validity-undecidable}.
\end{proof}
\end{annexitem}

\begin{annexitem}{annex:p04-hierarchy.2}{thm:p04-pairing-graph}{Pairing graph}
\begin{proof}
The atomic formula says in \(\mathcal N_{\mathrm{ar}}\):
\[
 (a+b)(a+b+1)+2b=2c,
\]
hence \(c=\operatorname{pair}_{\N}(a,b)\); doubling avoids division outside
\(\Lang_A\).  The bound \(b\leq_{\N}c\) is immediate, while
\(a\leq_{\N}a+b\leq_{\N}(a+b)(a+b+1)/2\leq_{\N}c\): the zero case forces \(a=0\), and
otherwise the relevant factor is at least \(1\).
\end{proof}
\end{annexitem}

\begin{annexitem}{annex:p04-hierarchy.3}{thm:p04-goedel-beta}{G\"odel beta function}
\begin{proof}
All quantifiers in \(\mathsf B\) are marked bounded and the remaining clauses
are atomic.  For fixed \(w\), pairing bijectivity gives unique \(t,m\);
monotonicity gives \(t,m\leq_{\N}w\).  Since \(d=1+(i+1)m\geq1\), Euclidean
division gives unique \(q,v\) with \(t=dq+v\), \(v<_{\N}d\),
\(q\leq_{\N}t<_{\N}St\), and \(v\leq_{\N}t\leq_{\N}w\).

Given \(a_0,\ldots,a_{k-1}\), set
\(N=\max(k,a_0,\ldots,a_{k-1})+1\) and \(m=N!\).  The numbers
\(1+(i+1)m\) for \(i<_{\N}k\) are pairwise coprime: a common prime divisor would
divide \((j-i)m\) but not \(m\), hence \(0<_{\N}j-i<_{\N}N\), and then divide
\(N!=m\), a contradiction.  Since \(a_i<_{\N}1+(i+1)m\), the Chinese remainder
theorem supplies \(t\) with those remainders.  Taking
\(w=\operatorname{pair}_{\N}(t,m)\) gives \(\beta(w,i)=a_i\) for every
\(i<_{\N}k\).
\end{proof}
\end{annexitem}

\begin{annexitem}{annex:p04-hierarchy.4}{thm:p04-pairing-contraction}{Pairing contraction}
\begin{proof}
Given witnesses \(a,b\), take \(c=\operatorname{pair}_{\N}(a,b)\).
The pairing lemma supplies \(\mathsf P(a,b,c)\) and \(a,b<_{\N}c+1\), proving the
forward implication; the reverse discards the pairing clause and bounds.
Freshness of \(z\) preserves \(\varphi\).  Apply the result to
\(\neg\varphi\) and negate to obtain the universal dual.
\end{proof}
\end{annexitem}

\begin{annexitem}{annex:p04-hierarchy.5}{thm:p04-bounded-exchange}{Bounded-quantifier exchange}
\begin{proof}
Fix an assignment and let \(n\) be the value of \(t\).  Right-to-left is
immediate.  Conversely, choose \(x_b\) for each \(b<_{\N}n\), and take \(c=0\)
if \(n=0\), otherwise \(c=1+\max_{b<_{\N}n}x_b\).  Then \(w=c\) bounds every
witness.  The side conditions preserve \(t\) and prevent capture.  Apply this
equivalence to \(\neg\varphi\) and negate to obtain the dual.
\end{proof}
\end{annexitem}

\begin{annexitem}{annex:p04-hierarchy.6}{thm:p04-prefix-normal-form}{Prefix normal form}
\begin{proof}
First, a prefix formula of level \(m<_{\N}n\) may be padded to either polarity at
level \(n\) by alternating vacuous quantifiers on fresh variables; \(\N\) is
nonempty.  Also write every implication introduced below as
\(\neg\mathsf P\lor\varphi\), preserving negation normal form.

Prove the closure assertion by simultaneous induction on \(n\), with the
\(\Pi\)-case dual.  At level \(0\), conjunction, disjunction, and marked
bounded quantifiers remain bounded syntactically.  For \(n\geq1\), repeated
pairing contraction converts a prefix \(\Sigma^0_n\)-formula to
\(\exists x\,\varphi\), where \(\varphi\) is prefix
\(\Pi^0_{n-1}\); the new bounded
quantifiers and pairing atoms are absorbed by the induction hypothesis.

Consider \(\exists x\varphi\) and \(\exists x'\psi\), with bound variables renamed
apart.  For conjunction and disjunction, pull both existential quantifiers
out, contract them to one paired witness, and use the induction hypothesis on
the resulting bounded matrix \(\varphi\land\psi\), respectively
\(\varphi\lor\psi\).
For marked bounded existence, first rename \(x\) so that \(x\ne y\) and
\(x\notin\operatorname{Var}(t)\), then commute
\(\exists y<t\exists x\varphi\) to \(\exists x\exists y<t\varphi\) and absorb the
bounded quantifier into \(\varphi\).  For marked bounded universality, bounded
exchange gives
\[
 \forall y<t\exists x\varphi
 \quad\Leftrightarrow\quad
 \exists w\forall y<t\exists x<w\varphi;
\]
absorb the two bounded quantifiers into \(\varphi\), then prepend \(\exists w\).
An additional unbounded existential is contracted with the existing first
block.  These are all generating operations.  Dualizing gives the
\(\Pi^0_n\) closure statement.

Now induct on the recursive generation of the syntactic classes.  Level zero
already is prefix form.  At level (n+1), an inherited level-(n) formula
first converts by induction and then pads; every closure operation converts
by the preceding paragraph.  The fixed negation-normal-form conversion
preserves satisfaction, so arbitrary formulas follow.  Conversely, every
prefix formula is generated by the syntactic clauses.  Hence the same
formulas define the same boldface relation classes, proving all three parts.
\end{proof}
\end{annexitem}

\begin{annexitem}{annex:p04-hierarchy.7}{thm:p04-bounded-verification}{Bounded verification of computations}
\sketchnote{Omitted: one displayed conjunction for (h1)--(h4) and the
digitwise check that each successor configuration follows its decoded
transition row.}
Translate the history clauses into conditions on entries of \(s\).  Extract
entries by bounded \(\mathsf B\), whose witnesses are bounded by \(s\), and
pairing components by \(\mathsf P\) with its monotonicity bounds.  Tape tag
\(0\) is the everywhere-blank tape; a positive tag decodes signed endpoints
and a fixed-length base-coded block, covering every finite-support tape.  Use
marked-bounded comparisons with fresh witnesses and numeral \(\bar2\) in the
binary recurrence.  Hence \(\mathsf{Comp}\) is bounded and states exactly the
history conditions.

If \(\varphi_e(x)\) is defined, its accepting computation, decoded table, and
certified digits form a finite sequence, which the beta lemma codes by some
\(s\).  Conversely, a satisfying \(s\) decodes a valid accepting computation
and output; invalid codes fail the designated-slot checks, matching everywhere
divergence.
\end{annexitem}

\begin{annexitem}{annex:p04-hierarchy.8}{thm:p04-ce-sigma-one}{Computable enumerability and level one}
\sketchnote{History coding inside \(\mathsf{Comp}\) is omitted.}
For c.e. \(A\), Theorem~\ref{thm:p04-enumeration} gives
\(A=\operatorname{dom}(\varphi_e)\), hence
\(x\in A\Leftrightarrow\exists s\,\mathsf{Comp}(\bar e,x,s)\), a
\(\Sigma^0_1\) definition.  Conversely search bounded witnesses.
\(\Delta^0_1\) means \(A\) and its complement are c.e.; the complement
criterion yields decidability.
\end{annexitem}
